% YY12100B
% Einschätzungsblock
% 14.0
% 2013-11-30
\label{m:ersteinschaetzung}\label{m:einschaetzungsblock}%
\par\noindent
% \TableWide
% 	{<Titel>}%1
% 	{<Beschreibung>}%2
% 	{<Label>}%3
% 	{<Quelle>}%4
% 	{<Lizenz>}%5
% 	{<Tabelle>}%6
% 	{<Float>}%8
\TableWide{Einschätzungsblock}
{}
{tab:einschaetzungsblock}
{}
{}
{
\CorrectionItemizeBox\renewcommand{\itemhook}{%
%   \setlength{\topsep}{0pt}%
  \setlength{\itemsep}{0.004\baselineskip}
%   \setlength{\parskip}{\baselineskip}
  \setlength{\leftmargin}{1.4em}
  \setlength{\labelwidth}{0.9em}
}
\noindent\begin{minipage}[t]{\linewidth -
{2\fboxrule} - {2\fboxsep}}
\noindent\begin{minipage}[t]{\linewidth}
{\vspace{-2\baselineskip}%
\scriptsize
% \begin{tabulary}{\linewidth}[H]{cp{8em}p{8.5cm}}
\begin{tabularx}{\linewidth}[H]{p{1.5cm}XX}
% \toprule\addlinespace
&
\TabHeading{Beurteilung / Untersuchung}
&
\TabHeading{Typische Sofortmaßnahmen}
\\\addlinespace\toprule\addlinespace
\noindent\TabHeading{Szeneüberblick m. Schutz}
&
\begin{minipage}[t]{\linewidth}
\begin{itemize}
  \item Sicherheit, Selbstschutz; Patientenanzahl; 
  \item Umgebung, Ort,
Zeit, Gefahrenzonen, Umstände,
\item Weitere Resourcen erforderlich? Unfallmechanismus,
Großschaden
\item Lagemedlung erforderlich
\item \EE{Trauma?} \EE{Mechanismus?}
\end{itemize}
\end{minipage}
&
\noindent\color{Indigo}\begin{minipage}[t]{\linewidth}
\begin{itemize}
  \item GAS-Maßnahmen 
  \item Lagemeldung
  \item Schutzausrüstung
  \item Nachforderung weiterer Kräfte
\end{itemize}
\end{minipage}
\\\addlinespace\toprule\addlinespace
\noindent\TabHeading{Eindruck}
&
\begin{minipage}[t]{\linewidth}
\begin{itemize}
  \item Alter, Geschlecht, \EE{Hautfarbe}, Gesichtsausdruck,
  \EE{Haltung}, spontane Bewegungen,  Sprache
  \item Offensichtliche Verletzungen, Blutungen
  \item Sonstige \EE{Auffälligkeiten}
\end{itemize}
\end{minipage}
&
\noindent\color{Indigo}\begin{minipage}[t]{\linewidth}
\begin{itemize}
  \item Manuelle HWS-Fixierung
  \item Stillung von starken Blutungen
  \item Bewegungsverbot
\end{itemize}
\end{minipage}
\\\addlinespace\midrule\addlinespace
\noindent\TabHeading{Bewusstsein}
&
\begin{minipage}[t]{\linewidth}
\begin{itemize}
  \item \EE{Bewusstseinsgrad} (WASB)
\end{itemize}
\end{minipage}
&
\noindent\color{Indigo}\begin{minipage}[t]{\linewidth}
\begin{itemize}
  \item NA-Nachforderung
\end{itemize}
\end{minipage}
\\\addlinespace\midrule\addlinespace
\noindent\TabHeading{Hauptbeschwerde}
&
\begin{minipage}[t]{\linewidth}
\begin{itemize}
  \item \EE{Berufungsgrund} und
  \EE{Leitsymptome}
\end{itemize}
\end{minipage}
&
\noindent\color{Indigo}\begin{minipage}[t]{\linewidth}
\end{minipage}
\\\addlinespace\toprule\addlinespace
\noindent\T{A}

\TabHeading{Atemweg}
&
\begin{minipage}[t]{\linewidth}
\begin{itemize}
  \item *Inspektion der Atemwege (\EE{Mund}, \EE{Nase})
  \item \EE{Atemgeräusche}
\end{itemize}
\end{minipage}
&
\noindent\color{Indigo}\begin{minipage}[t]{\linewidth}
\begin{itemize}
  \item Absaugung, Fremdkörper entfernen
  \item Kopf überstrecken, Esmarch-Handgriff
  \item HWS-Immobilisation manuell oder mit Schiene
  \item \Z{Erweiterte Maßnahmen}
\end{itemize}
\end{minipage}
\\\addlinespace\midrule\addlinespace
\noindent\T{B}

\TabHeading{Atmung}
&
\begin{minipage}[t]{\linewidth}
\begin{itemize}
  \item Schätzen: \EE{Atemfrequenz}, \EE{-tiefe}
  \item *\EE{Sp\ce{O2}}
  \item Atemgeräusche
  \item Inspektion der \EE{Thorax-Atembewegungen}, ggfs. 
  Atemprobe
  \item \E{Zeichen der
  Atemnot}: Hautfarbe, Bauch-/Thorax\-bewegungen,
  Anstrengung beim Atmen, Atemhilfsmuskulatur, \ldots
  \item \Z{(Auskultation: Seitenvergleich, Lungenbasen)}
\end{itemize}
\end{minipage}
&
\noindent\color{Indigo}\begin{minipage}[t]{\linewidth}
\begin{itemize}
  \item Reanimation
  \item Assistierte Beatmung (bei AF \textless\,8 oder
  \textgreater\,30\PerMin*)
  \item \ce{O2}-Gabe
  \item Situationsgerechte Lagerung
  \item \Z{Erweiterte Maßnahmen\\(Entlastung
  Spannungspneumothorax, \ldots )}
\end{itemize}
\end{minipage}
\\\addlinespace\midrule\addlinespace
\noindent\T{C}

\TabHeading{Kreislauf}
&
\begin{minipage}[t]{\linewidth}
\begin{itemize}
  \item \EE{Hautfarbe}, \EE{-temperatur}
  \item Schweiß
  \item \EE{Rekap}-Zeit
  \item \EE{Radialispuls}: Stärke, Rhythmus, Abschätzen der
  Frequenz
\end{itemize}
\end{minipage}
&
\noindent\color{Indigo}\itshape\begin{minipage}[t]{\linewidth}
\begin{itemize}
  \item Blutstillung
  \item Situationsgerechte Lagerung
\end{itemize}
\end{minipage}
\\\addlinespace\cmidrule{2-3}\addlinespace
\noindent\TabHeading{*STU}
&
\begin{minipage}[t]{\linewidth}
\begin{itemize}
  \item Inspektion
  und Abtasten von
  \begin{inparaenum}
    \item \EE{Kopf} inkl. Ohren,
    \item \EE{Hals},
    \item \EE{Thorax},
    \item \EE{Bauch},
    \item \EE{Becken} und
    \item \EE{Oberschenkel}
  \end{inparaenum}
\end{itemize}
\end{minipage}
&
\noindent\color{Indigo}\begin{minipage}[t]{\linewidth}
\begin{itemize}
  \item Blutstillung
  \item Situationsgerechte Lagerung
\end{itemize}
\end{minipage}
\\\addlinespace\midrule\addlinespace
\noindent\T{D}

\TabHeading{Schnelle\newline Untersuchung}
&
\begin{minipage}[t]{\linewidth}
\begin{itemize}
  \item Vitalwerte (\EE{HF}, \EE{RR})
  \item Neurologisch:
  \begin{inparaitem}
    \item \EE{Orientierung} oder *\EE{GCS}
    \item \EE{Pupillen}
    \item \EE{Kraftprobe OE}, *UE
    \item *Kann der
    Patient \EE{Hände} und \EE{Füße} spüren und aktiv
    bewegen?
  \end{inparaitem}
  \item *\EE{Blutzuckermessung}
\end{itemize}
\end{minipage}
&
\noindent\color{Indigo}\begin{minipage}[t]{\linewidth}
\begin{itemize}
  \item Transportentscheidung (eilig, Rendez-vous, \ldots)
\end{itemize}
\end{minipage}
\\\addlinespace\midrule\addlinespace
\noindent\T{E}

\TabHeading{Erweiterte\newline Untersuchung}
&
\begin{minipage}[t]{\linewidth}
\begin{itemize}
  \item (Fremd-) \EE{Anamnese} (SAMPLE, OPQRST)
  \item Einschätzen der \EE{Umgebung}
  \item Weitere relevante \EE{Untersuchungen}
  \item Arbeitsdiagnose erstellen und Differentialdiagnosen
  ausschließen
\end{itemize}
\end{minipage}
&
\noindent\color{Indigo}\begin{minipage}[t]{\linewidth}
\begin{itemize}
  \item Transportentscheidung 
  \item Spezielle Maßnahmen gemäß Verdachtsdiagnose
\end{itemize}
\end{minipage}
% \\\addlinespace\bottomrule\addlinespace
% \TabHeading{Beurteilung}
% &
% \begin{minipage}[t]{\linewidth}
% \begin{itemize}
%   \item \EE{Einschätzen der vitalen Bedrohung}
%   ständig während des gesamten Blocks
%   \item Vorläufige Verdachtsdiagnose formulieren
% \end{itemize}
% \end{minipage}
% &
% \noindent\color{Indigo}\begin{minipage}[t]{\linewidth}
% \begin{itemize}
%   \item {\TxMassVit}
%   \item Reihenfolge der weiteren Untersuchungen und
%   Maßnahmen festlegen
% \end{itemize}
% \end{minipage}
\\\addlinespace\bottomrule
\end{tabularx}

\vspace{-0.5\baselineskip}\begin{center}
  \EE{Der Einschätzungsblock muss in regelmäßigen Abständen
in angemessenem Umfang wiederholt
werden (Verlaufskontrolle)!}

{\tiny Mit einem * versehene Punkte werden nur durchgeführt,
wenn sie \E{situationsentsprechend} sind. \enquote{Typische
Sofortmaßnahmen} sind als \E{Beispiele} zu verstehen.
Bei \Ca{absolut zeitkritischen Patienten}, bei denen bereits die 
ABC-Einschätzung ergibt dass ein Transport nicht aufschiebbar ist, 
kann u.\,U. schon \Ca{nach C eine vorgezogene Strategie- 
und Transportentscheidung} getroffen werden. 
D und E sollen dann, sofern möglich, während des 
Transportes erfolgen. }
\end{center}
}
\end{minipage}
\end{minipage}\vspace{-\baselineskip}
}
{}

\vfill
\restoregeometry\pagestyle{fancy}

\Layout{57}{}
{}
{}
{%
\begin{minipage}{\linewidth}
\FontTableBase\footnotesize
{\sffamily\FontTableBase\footnotesize\CorrectionItemizeBox\renewcommand{\itemhook}{%
%   \setlength{\topsep}{0pt}%
  \setlength{\itemsep}{0.005\baselineskip}
%   \setlength{\parskip}{\baselineskip}
  \setlength{\leftmargin}{1.4em}
  \setlength{\labelwidth}{0.9em}
}
\begin{tabulary}{\linewidth}[H]{lCL}
\toprule\addlinespace
{\multirow{6}{2cm}{\TabHeading{Anamnese}}}
 &
{\sffamily\textbf S}
 &
Symptome {\&} Schmerzen: OPQRST

{\tiny(Lokalisation, Stärke,  Qualität, Zeit/Verlauf, 
 Ausstrahlung, Linderung/Verstärkung)}
\\\addlinespace
 &
{\sffamily\textbf A}
 &
Allergien/Allergische Reaktion (auch
Allergien auf Medikamente)
\\\addlinespace
 &
 {\sffamily\textbf M}
 &
Medikamente {\tiny(Nimmt Patient Medikamente?
Wogegen sind diese?)}
\\\addlinespace
 &
{\sffamily\textbf P}
 &
Patientengeschichte: Krankheiten {\tiny(chronische, frühere
K.)}, Operationen, \ldots 
\\\addlinespace
 &
 {\sffamily\textbf L}
 &
Letzte ... (zB letzte Mahlzeit, letzte
Regelblutung, letzter Spitalsaufenthalt)
\\\addlinespace
 &
 {\sffamily\textbf E}
 &
Ereignisse vor Notfalleintritt {\tiny(zB
Anstrengung vor Brustschmerz, Sturz, ...)}
\\\addlinespace\midrule\addlinespace

{\multirow{8}{2cm}{\TabHeading{Diagnostik}}}
 &
{\sffamily AF}
 &
Atemfrequenz auszählen
\\\addlinespace
 &
({\sffamily Sp\ce{O2}})
 &
Sauerstoffsättigung im Blut, Monitoring
\\\addlinespace
 &
{\sffamily  BZ}
 &
Blutzucker {\tiny(Wann war die letzte
Mahlzeit/Insulinspritze? Diabetes bekannt?)}
\\\addlinespace
 &
{\sffamily Temp.}
 &
Körpertemperatur/Fieber
\\\addlinespace
 &
 {\sffamily\small Neuro\-status}
 &
(Bewusstseinsgrad), Orientierung,
Pupillen, Halbseitenzeichen, Herdblick, Meningismus,
retrograde Amnesie
\\\addlinespace
 &
{\sffamily\small Trauma\-check}
 &
Abnorme Beweglichkeit, DMS an allen
Extremitäten (Nagelbettprobe), Wunden, paradoxe Atmung,
Abwehrspannung des Bauches, ...
\\\addlinespace
 &
{\sffamily\small Körp. Unt.}
&
Sonstige körperliche Untersuchungen: Abdomen abtasten,
Körperstellen inspizieren, \ldots
\\\addlinespace
 &
{\sffamily\small \ldots}
&
EKG, Kapnometrie, \ldots
\\\addlinespace\midrule\addlinespace
\TabHeading{Maßnahmen}
&
&
\begin{minipage}{\linewidth}
\begin{itemize}
  \item (Lagerung)
  \item Verband, Blutstillung
  \item Schienung
  \item psychischer Beistand
  \item \EE{Wärmeerhaltung / Kühlung}
  \item Entscheidung Notarzt-Nachforderung wegen
  Schmerztherapie oder Aufklärung/Belassung
  \item {\TxMassOGabe}
  \item Regelmäßige \EE{Wiederholung des
  Einschätzungsblockes} {\tiny(in angemessenem Umfang)}
  \item \EE{Immer durchzuführende Standardmaßnahmen}
  (\prettyref{m:standardmassnahmen-immer})
  \item \ldots
\end{itemize}
\end{minipage}
\\\addlinespace\midrule\addlinespace
\TabHeading{Assistenz}
&
&
\begin{minipage}{\linewidth}
\begin{itemize}
  \item EKG anlegen
  \item Infusion/periphervenösen Zugang vorbereiten
  \item Intubation vorbereiten
  \item Medikamente vorbereiten
  \item \ldots
\end{itemize}
\end{minipage}
\\\addlinespace\midrule\addlinespace
\TabHeading{Transport}
&
&
\begin{minipage}{\linewidth}
\begin{itemize}
  \item Welches Spital/Bett buche ich ab? {\tiny(Arbeitsunfall?
  Welche Abteilung fahre ich an? Ist der Patient bereits in
  einem Spital in Behandlung? Spitalsbehandlung vor kurzer
  Zeit (\enquote{Reklamation})?)}
  
  Brauche ich eine Spezialabteilung? {\tiny(Schockraum, Stroke
  Unit, PTCA, Verbrennungsstation)}
  \item Übergabe {\tiny(Anamnese, Befunde, Maßnahmen, Hinweise)}
\end{itemize}
\end{minipage}

\\\addlinespace\bottomrule
\end{tabulary}}

{\scriptsize\EE{Fett} gedruckte Punkte haben eine besonders hohe Priorität, (eingeklammerte) Punkte sollten schon als Teil des Einschätzungsblockes abgehandelt worden sein.}
\end{minipage}
}
{Kurzübersicht: \EEE{Weiterführende Untersuchungen und Anamnese}.
Die Reihenfolge ist je nach Patient
unterschiedlich!\label{tab:checkliste-patientenmanagement}\label{tab:pam-weiterfuehrende}.}
{}
